% ---------------------------------------------------------------------
\section{$R$ 比与流关联函数}\label{sec:ee-R}

在 $\sqrt s\ll M_Z$ 的能区，$e^+e^-$ 湮灭由单个虚光子传递；轻子一侧是 QED，强子一侧只经由电磁流的矩阵元出现，对末态求和后截面只依赖于一个标量函数。

\begin{definition}[$R$ 比]\label{def:ee-R}
在单光子交换、电磁耦合 $\alpha$ 的最低阶并忽略轻子质量时，定义
\begin{equation}\label{eq:ee-Rdef}
R(s)\equiv\frac{\sigma(e^+e^-\to\text{强子})}{\sigma(e^+e^-\to\mu^+\mu^-)},\qquad
\sigma(e^+e^-\to\mu^+\mu^-)=\frac{4\pi\alpha^2}{3s},
\end{equation}
$s=(k_1+k_2)^2$，$k_{1,2}$ 是 $e^\mp$ 的动量。记 $R_0\equiv N_c\sum_qe_q^2$，求和遍及该能区可以产生的夸克味，$e_q$ 以质子电荷为单位；$n_f=3,4,5$ 时 $R_0=2,\ \frac{10}3,\ \frac{11}3$。% VF-EE-1
% uses: none
\end{definition}

电磁流 $j^\mu=\sum_qe_q\bar q\gamma^\mu q$ 守恒且规范不变。按\cref{tab:front-notation}的约定定义% VF-CONV-12
\begin{align}
\Pi^{\mu\nu}(q)&\equiv\ii\!\int\!\dd^4x\,\ee^{\ii q\cdot x}\langle0|T\,j^\mu(x)j^\nu(0)|0\rangle=(q^\mu q^\nu-q^2g^{\mu\nu})\,\Pi(q^2), \label{eq:ee-Pidef}\\
W^{\mu\nu}(q)&\equiv\int\!\dd^4x\,\ee^{\ii q\cdot x}\langle0|j^\mu(x)j^\nu(0)|0\rangle
=\sum_X(2\pi)^4\delta^4(q-p_X)\,\langle0|j^\mu(0)|X\rangle\langle X|j^\nu(0)|0\rangle . \label{eq:ee-Wdef}
\end{align}
$\Pi^{\mu\nu}$ 的张量结构来自 Lorentz 协变性与流守恒（至多相差接触项，见\allowbreak\cref{sec:ee-analytic}）。\cref{prop:prelim-optical}（\tagR，$R_{\mathrm{def}}$）给出 $W^{\mu\nu}=2\Imag\Pi^{\mu\nu}$（$q^0>0$）；$W^{\mu\nu}$ 在横向子空间中半正定，故 $\Imag\Pi(s+\ii0)\ge0$。与 Peskin--Schroeder 的真空极化（$\Pi^{\mu\nu}_{\rm PS}=(q^2g^{\mu\nu}-q^\mu q^\nu)\Pi_{\rm PS}$）相比，在 $e$ 的最低阶 $\Pi_{\rm PS}=-e^2\Pi$。

\begin{proposition}[$R=12\pi\,\Imag\Pi$ \tagR]\label{prop:ee-R-ImPi}
（$R_{\mathrm{def}}$：由算符定义直接推出。）在\cref{def:ee-R}的条件下，
\begin{equation}\label{eq:ee-RImPi}
\sigma(e^+e^-\to\text{强子})=\frac{16\pi^2\alpha^2}{s}\,\Imag\Pi(s+\ii0),\qquad R(s)=12\pi\,\Imag\Pi(s+\ii0).
\end{equation}
% VF-CONV-12
非微扰地位：\cref{asm:intro-wightman}（MA0：末态的完备性与谱条件）。
\uses{def:ee-R,prop:prelim-optical,asm:intro-wightman}
\end{proposition}
\begin{proof}
在 $\alpha$ 的最低阶，振幅为 $\mathcal M=(e^2/s)\,\bar v(k_2)\gamma_\mu u(k_1)\langle X|j^\mu(0)|0\rangle$（光子传播子 $-\ii g_{\mu\nu}/q^2$，$q=k_1+k_2$；整体相位无关紧要）。轻子自旋平均给出 $L_{\mu\nu}\equiv\frac14\TrD[\sla k_2\gamma_\mu\sla k_1\gamma_\nu]=k_{1\mu}k_{2\nu}+k_{2\mu}k_{1\nu}-g_{\mu\nu}k_1\!\cdot\!k_2$。乘以通量因子 $1/(2s)$ 并对 $X$ 求和，由\cref{eq:ee-Wdef}与 $W^{\nu\mu}=2\Imag\Pi^{\nu\mu}$（\cref{prop:prelim-optical}），
\[
\sigma=\frac{e^4}{2s^3}\,L_{\mu\nu}W^{\nu\mu}(q)=\frac{e^4}{s^3}\,\Imag\Pi(s+\ii0)\,L_{\mu\nu}\big(q^\mu q^\nu-q^2g^{\mu\nu}\big).
\]
由 $k_i^2=0$、$k_1\cdot k_2=s/2$：$L_{\mu\nu}q^\mu q^\nu=2(k_1\cdot q)(k_2\cdot q)-s\,k_1\cdot k_2=0$（轻子流守恒），$L_{\mu\nu}g^{\mu\nu}=-s$，故 $L_{\mu\nu}(q^\mu q^\nu-q^2g^{\mu\nu})=s^2$，$\sigma=e^4\Imag\Pi/s=16\pi^2\alpha^2\Imag\Pi/s$。对单个点状轻子（$N_c=1$，电荷 $1$），\cref{eq:ee-ImPi0}给出 $\Imag\Pi_\mu=1/(12\pi)$，即 $\sigma_{\mu\mu}=4\pi\alpha^2/(3s)$；相除得 $R=12\pi\Imag\Pi$。对强子末态求和用到末态的完备性，即非微扰地位一行所列的假设。
\end{proof}
对自由无质量夸克，单圈关联函数给出
\begin{equation}\label{eq:ee-ImPi0}
\Imag\Pi_0(s+\ii0)=\frac{N_c}{12\pi}\sum_qe_q^2,\qquad R=R_0,\qquad \sigma_0=\frac{4\pi\alpha^2}{3s}R_0 ,
\end{equation}
% VF-EE-1, VF-EE-3
$D$ 维中的完整计算（$\overline{\rm MS}$ 形式 $\Pi_0=\frac{N_ce_q^2}{12\pi^2}(\frac53-L)$）见\cref{ex:prelim-vacpol}。

\paragraph{$D$ 维的记账。}\cref{sec:ee-oneloop}在 $D=4-2\varepsilon$ 维中分别计算实辐射与虚修正，那时只比较强子一侧的量。定义
\begin{equation}\label{eq:ee-rho}
\rho(s)\equiv-g_{\mu\nu}W^{\mu\nu}(q)=2(D-1)\,s\,\Imag\Pi(s+\ii0),
\end{equation}
第二个等号在任意 $D$ 中成立，因为 $-g_{\mu\nu}(q^\mu q^\nu-q^2g^{\mu\nu})=(D-1)q^2$。领头阶 $\rho_0=N_ce_q^2\,2(D-2)s\,\Phi_2$，其中 $-g_{\mu\nu}\TrD[\sla p_1\gamma^\mu\sla p_2\gamma^\nu]=2(D-2)s$，$\Phi_2$ 是 $D$ 维两体相空间（\cref{lem:appC-phase-space}）。由于各阶的极点在 $\rho$ 的和中相消（\cref{thm:ee-R-finite}），
\begin{equation}\label{eq:ee-Rratio}
\frac{R^{(n)}}{R_0}=\lim_{\varepsilon\to0}\frac{\rho^{(n)}}{\rho_0}
\end{equation}
与轻子一侧如何延拓到 $D$ 维无关。
